\documentclass[10pt,a4paper]{article}
\usepackage[utf8]{inputenc}
\usepackage[T1]{fontenc}
\usepackage[a4paper,top=20mm,bottom=20mm,left=16mm,right=16mm,headheight=14pt]{geometry}
\usepackage{fancyhdr}
\usepackage{booktabs}
\usepackage{longtable}
\usepackage{array}
\usepackage{enumitem}
\usepackage{amssymb}
\usepackage{xcolor}
\usepackage[hyphens]{url}
\usepackage[colorlinks=true,urlcolor=blue!60!black,linkcolor=blue!60!black]{hyperref}

\setlength{\parindent}{0pt}
\setlength{\parskip}{3pt}
\renewcommand{\arraystretch}{1.28}

\pagestyle{fancy}
\fancyhf{}
\fancyhead[L]{\small LANGUAGE \& SOCIETY}
\fancyhead[R]{\small TRANSCRIPT ANALYSIS WORKSHEET}
\fancyfoot[C]{\small\thepage}
\renewcommand{\headrulewidth}{0.4pt}

\begin{document}

\begin{center}
    {\large\bfseries LANGUAGE \& SOCIETY PROJECT}\\[3pt]
    {\LARGE\bfseries Transcript Analysis Worksheet}\\[3pt]
    {\normalsize Qualitative analysis of political news discourse}
\end{center}
\vspace{2pt}
\noindent\rule{\textwidth}{0.4pt}
\vspace{4pt}


\section*{Creator 1: Sarthak Goswami (\textit{Sunday Show} --- Independent Satirical Infotainment)}
\begin{itemize}[leftmargin=1.5em, itemsep=3pt, topsep=4pt]
  \item \textbf{Analyst:} Baani Solanki (\texttt{2025114016})
  \item \textbf{Channel / Outlet:} Sarthak Goswami (\texttt{@SundaySarthak}, YouTube)
  \item \textbf{Ideological / Editorial Positioning:} Centrist-liberal satirical infotainment; critiques ruling-party majoritarianism ("dharma ka dhandha"), opposition opportunism ("sasti BJP"), bureaucratic corruption, and mainstream "godi media" polarization while foregrounding civic constitutionalism and youth agency.
  \item \textbf{Transcript Files Referenced:}
  \begin{itemize}[leftmargin=1.4em, itemsep=2pt, topsep=2pt]
    \item Event 1: \texttt{transcripts/event1\_ram\_mandir/01\_sarthak\_goswami\_lb7gZnyZH0c.txt}
    \item Event 2: \texttt{transcripts/event2\_neet\_protest/01\_sarthak\_goswami\_EM1GeZufO6E.txt}
  \end{itemize}
\end{itemize}
\vspace{4pt}\noindent\rule{\textwidth}{0.4pt}\vspace{4pt}

\section*{Guidelines: Five features to look for in every event}
\begin{longtable}{@{}>{\raggedright\arraybackslash}p{0.17\textwidth} >{\raggedright\arraybackslash}p{0.20\textwidth} >{\raggedright\arraybackslash}p{0.24\textwidth} >{\raggedright\arraybackslash}p{0.32\textwidth}@{}}
\toprule
\textbf{Feature} & \textbf{What to look for} & \textbf{Illustrative example} & \textbf{Questions to answer} \\
\midrule
\endhead
\textbf{1. Naming} & What are people, groups, actions, or events called? & "Students", "protesters", "agitators"; "irregularity", "scam", "controversy". & What exact label is used? Who or what does it refer to? Who uses it? What alternative label could have been used? \\
\addlinespace[4pt]
\textbf{2. Presupposition} & What is treated as already true or taken for granted? & "Why are students creating chaos again?" may presuppose earlier disruption. & What does the wording assume? Is that assumption supported, disputed, or left unexamined? Does it appear in a question or words such as "again", "still", or "finally"? \\
\addlinespace[4pt]
\textbf{3. Agency} & Who is shown doing something? Who is acted upon? Who disappears? & "Police injured protesters" vs. "Protesters were injured." & Who is the actor? Who experiences the action? Is responsibility named, backgrounded, or omitted? \\
\addlinespace[4pt]
\textbf{4. Antagonisation} & Who is positioned against whom? & "Ordinary citizens" vs. "political troublemakers". & Which groups are contrasted? Is one presented as more legitimate, moral, patriotic, dangerous, or disloyal? Is an "us vs. them" division implied? \\
\addlinespace[4pt]
\textbf{5. Intensification / softening} & What is exaggerated, minimised, softened, or euphemised? & "Violent mob" vs. "peaceful demonstrators"; "incident" instead of a direct description. & Which words heighten emotion or seriousness? Which minimise it? Are metaphors, dramatic adjectives, hedges, or euphemisms used? Does the wording match the evidence? \\
\bottomrule
\end{longtable}
\textit{Reminder: A label or implication is not automatically false or manipulative. Examine the surrounding context and evidence.}

\vspace{4pt}\noindent\rule{\textwidth}{0.4pt}\vspace{4pt}

\section*{Event1. Ram Mandir donation controversy}
\textbf{Focus:} How is the controversy framed---as financial accountability, religious sentiment, political motivation, or something else? Be precise about the specific allegation or episode covered.

\begin{itemize}[leftmargin=1.5em, itemsep=3pt, topsep=4pt]
  \item \textbf{Outlet:} Sarthak Goswami (\textit{The Sunday Show}, YouTube)
  \item \textbf{Date:} \texttt{2026-07-05}
  \item \textbf{Video title:} \textit{Ram Mandir Scam, Gadkari's Ethanol, Kejriwal's Temple \& More | Sunday Show} (\url{https://youtu.be/lb7gZnyZH0c?si=QvgBq0ACOhcx0pdO})
  \item \textbf{Timestamp(s):} Relevant Event 1 segment \texttt{[00:00--06:30]} (with brief callbacks at \texttt{[27:55--28:15]} and \texttt{[29:15--29:20]}; total video duration \texttt{45:41}; unrelated weekly news segments \texttt{[06:31--45:41]} and Odoo POS sponsor segment \texttt{[07:56--10:28]} excluded per project guidelines. Contextual follow-up in 26 July episode \texttt{EM1GeZufO6E} at \texttt{[27:17--29:45]} also noted).
  \item \textbf{Analyst:} Baani Solanki (\texttt{2025114016})
  \item \textbf{Language:} Hindi / Hinglish (Native Hindi audio, Romanized script; non-AI-dubbed)
\end{itemize}

\subsection*{Event-specific questions (Event 1)}
\begin{enumerate}[leftmargin=1.6em, itemsep=4pt, topsep=4pt]
  \item \textbf{What terms are used: "scam", "alleged irregularity", "corruption", "political controversy", or something else?}
  \begin{itemize}[leftmargin=1.4em, itemsep=2pt, topsep=2pt]
    \item Sarthak uses blunt vernacular terms of financial crime---\texttt{ghotala} [scam] (\texttt{[00:00]}, \texttt{[00:19]}, \texttt{[00:46]}), \texttt{paisa chori} / \texttt{chanda chori} [donation theft] (\texttt{[00:00]}, \texttt{[00:19]}, \texttt{[00:29]}), \texttt{ghapla} [bungling/fraud] (\texttt{[02:21]}), and \texttt{skaimon} [scams] (\texttt{[01:31]})---before reframing the controversy metaphorically as \texttt{asli ghotala} [the real scam] (\texttt{[00:10]}, \texttt{[04:04]}, \texttt{[04:30]}, \texttt{[05:22]}, \texttt{[06:16]}) and \texttt{Ram ka charitra chori} [theft of Ram's character/ethos] (\texttt{[06:26]}). He also coins the satirical political hybrid \texttt{"kangresiphikeshan of BJP politiks"} [Congressification of BJP politics] (\texttt{[00:29]}).
  \end{itemize}
  \item \textbf{How are the temple, trust, donors, and people making allegations described?}
  \begin{itemize}[leftmargin=1.4em, itemsep=2pt, topsep=2pt]
    \item \textbf{Trust \& Accused Staff:} Senior trustee Champat Rai is labelled \texttt{"RSS VHP ka aadmi"} [an RSS-VHP man] (\texttt{[00:46]}) who claims innocence by blaming his driver \texttt{"salman khaan se inspireshan le rahe hain"} [taking inspiration from Salman Khan's hit-and-run defense] (\texttt{[01:31--01:40]}). Driver Ram Shankar "Tinnu" Yadav is simultaneously cited as the alleged \texttt{"mastarmaind"} (\texttt{[00:37]}) and satirized as \texttt{"bali ka bakra"} [scapegoat] (\texttt{[00:46]}). Counting staff like Lavkush Mishra are mocked for treating temple donations as \texttt{"hamare baap ka maal"} [our father's property] (\texttt{[01:14]}).
    \item \textbf{Donors / Public:} Devotees are depicted with biting second-person irony---both as \texttt{"bhakt bechare kamphyooj ho gae hain"} [poor devotees are confused about how to accuse one 'god' of stealing from another God] (\texttt{[01:48--01:56]}) and as complicit citizens (\texttt{"aap bhi bhagidar ho usmen"}, \texttt{[02:29]}; \texttt{"dharma chalane ka bachane ka theka toh tum hi ne diya tha"}, \texttt{[02:46--02:54]}) who outsourced their religion to political contractors. Online hyper-partisans defending the theft are mocked for arguing \texttt{"Hindu ka paisa agar kisi Hindu ne le bhi liya... galat kya hai?"} (\texttt{[02:54--03:12]}).
    \item \textbf{People Making Allegations:} Local institutional actors like the Ayodhya \texttt{"baar esosieshan"} [Bar Association] threatening a city blockade (\texttt{[01:48]}) are cited alongside \texttt{"apojishan ne alag baval kaat diya hai"} [the opposition has raised its own ruckus] (\texttt{[01:48]}).
  \end{itemize}
  \item \textbf{Is the issue framed as financial accountability, an attack on Hindu faith, a political dispute, or a matter of public trust?}
  \begin{itemize}[leftmargin=1.4em, itemsep=2pt, topsep=2pt]
    \item Primarily framed as \textbf{moral-civic accountability and the commercialization/politicization of religion (\texttt{"dharma ka dhandha"}, \texttt{[02:21]})}. Rather than treating the controversy as an external attack on Hindu faith, Sarthak argues that Hindu faith was hollowed out \textit{from within} by self-styled religious gatekeepers (\texttt{"adharmi log jab dharma ka prachar karne lage wo tha asli ghotala"}, \texttt{[04:04]}).
  \end{itemize}
  \item \textbf{Are questions about donations presented as legitimate scrutiny or as an attempt to malign the temple or its supporters?}
  \begin{itemize}[leftmargin=1.4em, itemsep=2pt, topsep=2pt]
    \item Presented unequivocally as \textbf{legitimate, overdue scrutiny}. In fact, Sarthak ridicules partisan apologists who gaslight questioners (\texttt{"tumhare dharma ka theka leke tumko hi loot liya aur ab tum sawal karoge"}, \texttt{[03:12--03:20]}).
  \end{itemize}
  \item \textbf{Are allegations attributed to specific people and supported with evidence, or presented as established facts?}
  \begin{itemize}[leftmargin=1.4em, itemsep=2pt, topsep=2pt]
    \item Hybrid: Financial recoveries and arrests are anchored to the SIT investigation (\texttt{"SIT ki jaanch baithi jinhonne apni report dayar ki jiske baad Trust ke aath logon ko arest kiya gaya"}, \texttt{[00:29--00:37]}) with concrete figures (\texttt{Rs.\,20 lakh} at Avinash Shukla's yoga center, \texttt{Rs.\,12,000} salary vs \texttt{Rs.\,25 lakh} car and \texttt{Rs.\,12 lakh} stable for Lavkush Mishra at \texttt{[01:14--01:23]}). However, his broader claim that senior figures used the driver as a \texttt{"bali ka bakra"} (\texttt{[00:46]}) is presented through satirical inference (\texttt{"ab aap thoda Google karo pata lag jaega"}, \texttt{[00:46]}) rather than adjudicated proof.
  \end{itemize}
  \item \textbf{How are denials and counterarguments represented: explained, dismissed, challenged, or omitted?}
  \begin{itemize}[leftmargin=1.4em, itemsep=2pt, topsep=2pt]
    \item \textbf{Mocked and challenged through satire.} Champat Rai's denial (\texttt{"champat raay ji ka kahna hai ki ji main toh beksoor hoon. mere draivar ne mujhe dhokha diya hai"}, \texttt{[01:31--01:40]}) is immediately ridiculed via the Salman Khan driver-blame trope (\texttt{[01:40]}). Partisan online counterarguments ("if a Hindu took a Hindu's money, why are you upset?") are quoted and dismantled through reductio ad absurdum (\texttt{"kal koi goli maar dega tumhen aur yeh kahenge ki Hindu ka khoon Hindu ne bahaya"}, \texttt{[03:12]}).
  \end{itemize}
  \item \textbf{Does the presenter distinguish criticism of the trust or its financial practices from criticism of the temple or religion itself?}
  \begin{itemize}[leftmargin=1.4em, itemsep=2pt, topsep=2pt]
    \item \textbf{Yes, explicitly and sharply.} He contrasts the ancient tradition of Hinduism (\texttt{"hajaron lakhon saal purana dharma"}, \texttt{[03:38]}) and Lord Ram's moral character against opportunistic political contractors and online abusers (\texttt{"jo Mandir ke bahar chappal churate the unko andar ka theka de diya tumne... Ram ka chanda nahin Ram ka charitra chori ho chuka hai"}, \texttt{[06:16--06:26]}).
  \end{itemize}
  \item \textbf{Who is portrayed as protecting donors' interests, and who as exploiting religious sentiment?}
  \begin{itemize}[leftmargin=1.4em, itemsep=2pt, topsep=2pt]
    \item \textbf{Exploiting religious sentiment:} Trust insiders, VIP-darshan profiteers (\texttt{[02:29--02:38]}), BJP-RSS-linked contractors who peddled the \texttt{"Hindu khatre mein hai"} [Hindus are in danger] tonic (\texttt{[03:20--03:28]}), and abusive online "dharma rakhvala" trolls (\texttt{[04:30--05:22]}).
    \item \textbf{Protecting donors' interests:} No political party is idealized as a savior (later at \texttt{[26:44--27:55]} he equally mocks Kejriwal/AAP as \texttt{"sasti BJP"} for temple-appeasement politics); instead, Sarthak calls upon ordinary Hindu citizens themselves to wake up (\texttt{"Hindu jago"}, \texttt{[06:26]}) and reclaim ethical accountability.
  \end{itemize}
\end{enumerate}

\subsection*{Apply the five shared features (Event 1)}
\begin{longtable}{@{}>{\raggedright\arraybackslash}p{0.19\textwidth} >{\raggedright\arraybackslash}p{0.38\textwidth} >{\raggedright\arraybackslash}p{0.37\textwidth}@{}}
\toprule
\textbf{Feature} & \textbf{Exact wording + timestamp} & \textbf{What it does in context} \\
\midrule
\endhead
\textbf{1. Naming} & \texttt{"champat raay RSS VHP ka aadmi hai aur tinnoo draivar bali ka bakra."} \texttt{[00:46]} \newline  \texttt{"laik vhat ij dis kangresiphikeshan of BJP politiks."} \texttt{[00:29]} & Labels Champat Rai by his ideological patronage network (\texttt{RSS VHP ka aadmi}) while naming arrested driver Tinnu Yadav a \texttt{"bali ka bakra"} (scapegoat), challenging the official narrative that only low-level staff were culpable. \texttt{"Kangresiphikeshan"} equates BJP-era institutional graft with older Congress-era corruption scandals. \\
\addlinespace[4pt]
\textbf{2. Presupposition} & \texttt{"par tum mujhe batao tum sarapraizd kyon ho? ... jab itne salon se dharma ka dhandha ban gaya tha toh usmen ghapla toh hona hi tha."} \texttt{[02:12--02:21]} & The \textit{why}-question (\texttt{tum sarapraizd kyon ho?}) and temporal clause (\texttt{jab itne salon se dharma ka dhandha ban gaya tha}) presuppose as an established fact that religion had already been turned into a commercial racket (\texttt{dhandha}) long before the theft surfaced. \\
\addlinespace[4pt]
\textbf{3. Agency} & \texttt{"dharma chalane ka bachane ka theka toh tum hi ne diya tha in logon ko... aap hi ne nilam kiya aur ab aap boli ka hisab bhi mang rahe ho."} \texttt{[02:46--03:38]} & Shifts agency directly onto the viewer/voter (\texttt{tum hi ne}, \texttt{aap hi ne}) using active ergative constructions (\texttt{ne diya tha}, \texttt{ne nilam kiya}), refusing to let the public pose purely as passive victims of corrupt trustees. \\
\addlinespace[4pt]
\textbf{4. Antagonisation} & \texttt{"jo Mandir ke bahar chappal churate the unko andar ka theka de diya tumne... Ram ka chanda nahin Ram ka charitra chori ho chuka hai."} \texttt{[06:16--06:26]} & Constructs a sharp moral opposition between authentic ethical devotion (\texttt{Ram ka charitra}) and lumpenized, abusive self-appointed guardians of Hindutva (\texttt{chappal churate the}, online abusers chanting \texttt{"jay shri Ram"} while issuing rape threats at \texttt{[04:47--05:22]}). \\
\addlinespace[4pt]
\textbf{5. Intensification / softening} & \textbf{Irony/Hook reversal:} \texttt{"Ram Mandir mein ghotala nahin huaa hai... kyunki asli ghotala hi kahin aur huaa hai."} \texttt{[00:00--00:10]} \newline  \textbf{Intensification:} \texttt{"sher ko pala tha ab uski dait se dikkat... Hindu khatre mein tumhen ghutti pilae ja rahe the."} \texttt{[03:20--03:28]} & Opens with mock-denial (ironic softening) as a rhetorical hook, then intensifies critique through visceral metaphors---raising a predator (\texttt{sher ko pala}), force-feeding fear (\texttt{ghutti pilae}), and auctioning religion (\texttt{nilam kiya}, \texttt{[03:28]}). \\
\bottomrule
\end{longtable}

\subsection*{Short summary (Event 1)}
\begin{itemize}[leftmargin=1.5em, itemsep=3pt, topsep=4pt]
  \item \textbf{What interpretation of the event does this video encourage?} That the Ayodhya donation embezzlement is the inevitable outcome (\texttt{"prakriti ka niyam"}, \texttt{[02:21]}) of mixing electoral politics, unchecked VIP commercialization, and blind majoritarian loyalty (\texttt{[02:12--03:38]}), where junior staff like driver Tinnu Yadav take the fall (\texttt{[00:46]}) while deeper moral corruption---weaponizing Lord Ram's name for power and online abuse (\texttt{[04:30--06:26]})---goes unpunished.
  \item \textbf{Whose perspective is foregrounded, challenged, or missing?} \textit{Foregrounded:} Secular-liberal / self-reflective Hindu citizen critique and victims of online Hindutva harassment (\texttt{[04:38--05:48]}). \textit{Challenged:} Trust leadership's "rogue driver" excuse (\texttt{[01:31--01:40]}) and online bhakt apologia (\texttt{[02:54--03:12]}). \textit{Missing:} Detailed institutional explanation from SBI or SIT officials and the perspective of sincere grassroots temple volunteers.
  \item \textbf{What evidence supports your interpretation? What complicates it?} \textit{Supports:} He cites SIT arrests of 8 staffers (\texttt{[00:37]}), specific asset seizures (\texttt{Rs.\,20 lakh} from Avinash Shukla, Lavkush Mishra's \texttt{Rs.\,12,000} salary vs luxury assets at \texttt{[01:14--01:23]}), and verbatim screenshots/quotes of abusive online Hindutva accounts (\texttt{[04:47--05:22]}). \textit{Complicates:} Over half of the 6.5-minute segment (\texttt{[03:47--06:15]}) pivots away from the Ayodhya Trust's accounting mechanics to an unrelated viral online harassment incident to make a broader cultural point, blurring financial evidence with cultural commentary.
\end{itemize}
\vspace{4pt}\noindent\rule{\textwidth}{0.4pt}\vspace{4pt}

\section*{Event2. Students' protest regarding NEET paper leak}
\textbf{Focus:} How does coverage construct the legitimacy of student protest and responsibility for the underlying problem?

\begin{itemize}[leftmargin=1.5em, itemsep=3pt, topsep=4pt]
  \item \textbf{Outlet:} Sarthak Goswami (\textit{The Sunday Show}, YouTube)
  \item \textbf{Date:} \texttt{2026-07-26}
  \item \textbf{Video title:} \textit{Dharmendra Pradhan Resigned, CJP won, Modi on Instagram \& more | Sunday Show} (\url{https://youtu.be/EM1GeZufO6E?si=PgvQ7GHqNjZJIZET})
  \item \textbf{Timestamp(s):} Relevant Event 2 segments \texttt{[00:00--08:55]} and \texttt{[10:20--23:38]} (total video duration \texttt{33:04}; Kiwi sponsor ads \texttt{[08:56--10:20]}, \texttt{[20:52--21:03]} and unrelated news segments \texttt{[23:39--33:04]} excluded per project guidelines).
  \item \textbf{Analyst:} Baani Solanki (\texttt{2025114016})
  \item \textbf{Language:} Hindi / Hinglish (Native Hindi audio, Romanized script; non-AI-dubbed)
\end{itemize}

\subsection*{Event-specific questions (Event 2)}
\begin{enumerate}[leftmargin=1.6em, itemsep=4pt, topsep=4pt]
  \item \textbf{Are participants called "students", "candidates", "protesters", "agitators", "mobs", or "political pawns"?}
  \begin{itemize}[leftmargin=1.4em, itemsep=2pt, topsep=2pt]
    \item Sarthak consistently legitimizes participants as \texttt{"students"} (\texttt{[00:08]}), \texttt{"aam janta"} [ordinary public] (\texttt{[03:39]}), \texttt{"jenjis"} / Gen-Z youth (\texttt{[10:40]}, \texttt{[12:04]}), \texttt{"lakhon bachche"} [lakhs of children/kids] (\texttt{[22:48]}), and \texttt{"lidarles moovment"} whose leader became \texttt{"khud pablik"} [the public itself] (\texttt{[08:08]}). He explicitly quotes and mocks mainstream TV anchors who label them \texttt{"aatankvadi"} [terrorists] (\texttt{[05:31]}, \texttt{[22:22]}), \texttt{"goons oph da CJP"} (\texttt{[05:38]}), \texttt{"desh virodhi gaing"} [anti-national gang] (\texttt{[05:47]}), and \texttt{"pakistan phanded"} (\texttt{[22:48]}). He also honestly notes peripheral \texttt{"ansoshal eliments"} [anti-social elements] hooting near Connaught Place (\texttt{[07:12--07:21]}), separating them from the core Jantar Mantar site.
  \end{itemize}
  \item \textbf{Are their stated demands explained, summarised, mocked, questioned, or left out?}
  \begin{itemize}[leftmargin=1.4em, itemsep=2pt, topsep=2pt]
    \item \textbf{Clearly explained, summarised, and endorsed.} He opens with Education Minister Dharmendra Pradhan's resignation (\texttt{[00:08]}) and itemizes the remaining three accepted demands at \texttt{[08:16--08:35]}: (i) \texttt{Rs.\,1 crore} compensation to families of students who died by suicide due to paper leaks, (ii) a formal Delhi Police apology for the 20 July crackdown, and (iii) withdrawal/non-filing of criminal cases against protesters, alongside deeper systemic anti-leak reforms (\texttt{[16:25--17:07]}).
  \end{itemize}
  \item \textbf{Does the presenter focus on the original grievance, disruption, political affiliations, or public inconvenience?}
  \begin{itemize}[leftmargin=1.4em, itemsep=2pt, topsep=2pt]
    \item Focuses squarely on the \textbf{original systemic grievance} (repeated paper leaks, student suicides, collapsing education budget from \texttt{4.5\%} to \texttt{2.5\%} of GDP, \texttt{[15:59--18:10]}) and \textbf{democratic civic solidarity} (\texttt{[07:21--08:16]}), while debunking claims of foreign funding (\texttt{[00:40--00:48]}).
  \end{itemize}
  \item \textbf{Who is represented as responsible: examination authorities, government, students, political parties, or others?}
  \begin{itemize}[leftmargin=1.4em, itemsep=2pt, topsep=2pt]
    \item Primary responsibility is assigned to the \textbf{Union \& State Governments} (for reactive \texttt{"ditrens baay phiyar"} instead of preventive exam reforms and for slashing education spending in favor of electoral \texttt{"phri bis"} [freebies], \texttt{[16:34--18:53]}), \textbf{Delhi Police \& unbadged plainclothes men} for the 20 July crackdown (\texttt{[03:39--04:17]}), \textbf{Central Universities} (DU, JNU, Jamia, IITs) for issuing circulars threatening students with legal action (\texttt{[15:00--15:20]}), and \textbf{mainstream TV anchors} for vilifying youth (\texttt{[05:31--05:56]}, \texttt{[21:56--23:30]}).
  \end{itemize}
  \item \textbf{Are questions framed around whether students have legitimate grievances, or why they are causing disruption?}
  \begin{itemize}[leftmargin=1.4em, itemsep=2pt, topsep=2pt]
    \item Framed entirely around \textbf{students' legitimate grievances and state accountability}: why it took a massive street movement for a single resignation (\texttt{[01:51]}, \texttt{[15:59--16:16]}), who ordered unbadged men to assault female protesters (\texttt{[03:57--04:17]}), and why education budgets are shrinking (\texttt{[17:07--18:44]}).
  \end{itemize}
  \item \textbf{Are claims of violence, disorder, or political manipulation attributed to sources and substantiated?}
  \begin{itemize}[leftmargin=1.4em, itemsep=2pt, topsep=2pt]
    \item Claims of foreign/political conspiracy are refuted by citing the \textbf{Ministry of External Affairs (MEA)} clarification that there was \texttt{"koi bhi pakistan, chaina ya phoren phanding nahin"} (\texttt{[00:40--00:48]}). Claims of police violence (\texttt{lathi charj}, \texttt{aansoo gais}, \texttt{pelet gans}, plainclothes harassment at \texttt{[03:39--04:17]}) are backed by ground eyewitness observation (\texttt{[07:12--08:16]}) and cellphone videos (\texttt{[04:53]}), while fairly acknowledging that \texttt{"wahan pe kaee sare police vale bhi hart hue the. ritalieshan mein un pe bhi hamla huaa tha"} (\texttt{[04:45--04:53]}). He also transparently notes CJP founder Abhijit Dipke's prior AAP social-media stint (\texttt{[02:19]}).
  \end{itemize}
  \item \textbf{How are police and government actions described compared with protesters' actions? Are similar actions evaluated differently?}
  \begin{itemize}[leftmargin=1.4em, itemsep=2pt, topsep=2pt]
    \item Police actions on 18 July and 20 July are condemned in stark active verbs: \texttt{"abadakt karke le jati hai. not diten abadakt"} [abducts, not detains Sonam Wangchuk] (\texttt{[03:02--03:12]}), \texttt{"broot phors yooz karke"} (\texttt{[04:27]}), \texttt{"pelet gans ki bauchhar"} (\texttt{[03:39--03:47]}). Conversely, protester clashes are contextualized as \texttt{"ritalieshan"} [retaliation] (\texttt{[04:45]}), and post-20-July protesters are celebrated for peaceful mutual aid---forming human chains, medical camps, book libraries, and Zomato/Swiggy food sharing (\texttt{[07:21--08:08]}). Interestingly, at \texttt{[11:54--12:31]}, he humanizes young Gen-Z police jawans on the ground as sympathetic peers caught under senior orders.
  \end{itemize}
  \item \textbf{Are students allowed to explain their demands in their own words, or does the presenter mainly speak for them?}
  \begin{itemize}[leftmargin=1.4em, itemsep=2pt, topsep=2pt]
    \item The presenter \textbf{mainly speaks for and narrates on behalf of the students} in his studio monologue format, supplemented by reading out student protest posters (\texttt{"pavar of pipal ij gretar den da pipal in pavar"}, \texttt{[12:39]}) and recounting his own visit to Jantar Mantar (\texttt{[07:12--08:16]}), while inserting brief clips of politicians and rival TV anchors (\texttt{[05:38--05:56]}, \texttt{[22:30]}).
  \end{itemize}
  \item \textbf{Does the outlet's framing differ between coverage of the NEET paper-leak issue before the protest and coverage once students begin protesting?}
  \begin{itemize}[leftmargin=1.4em, itemsep=2pt, topsep=2pt]
    \item Rather than turning hostile once students hit the streets (as mainstream TV channels did), Sarthak argues that the street protest \textbf{elevated and validated} his prior anti-corruption stance on education---declaring \texttt{"bhaad mein gaee nyootraliti. agar tum ejukeshan karapt kar doge toh desh barbad hai"} (\texttt{[18:17--18:26]}). \textit{(Note: Direct pre-protest episodes are outside this single transcript, though he references his continuous week-long coverage at \texttt{[32:36]}.)}
  \end{itemize}
\end{enumerate}

\subsection*{Apply the five shared features (Event 2)}
\begin{longtable}{@{}>{\raggedright\arraybackslash}p{0.19\textwidth} >{\raggedright\arraybackslash}p{0.38\textwidth} >{\raggedright\arraybackslash}p{0.37\textwidth}@{}}
\toprule
\textbf{Feature} & \textbf{Exact wording + timestamp} & \textbf{What it does in context} \\
\midrule
\endhead
\textbf{1. Naming} & \texttt{"subah-subah hangar straik par baithe sonam vangachuk ko dilli police aati hai sivil kloths mein... wahan se abadakt karke le jati hai. not diten abadakt"} \texttt{[03:02--03:12]} & Explicitly rejects the legalistic police euphemism \texttt{"diten"} (detain) and replaces it with \texttt{"abadakt"} (abduct), criminalizing the plainclothes removal of hunger-striker Sonam Wangchuk and framing 20 July as \texttt{"blaik mande"} [Black Monday] (\texttt{[03:12]}). \\
\addlinespace[4pt]
\textbf{2. Presupposition} & \texttt{"aakhiri mein eego todna hi pada... pipal aar not skerd enimor."} \texttt{[00:40--00:58]} & \texttt{"Eego todna hi pada"} presupposes that the government's 30-day delay was driven by personal arrogance (\texttt{ego}) rather than policy deliberation; \texttt{"not skerd enimor"} presupposes an earlier climate of state-induced public fear. \\
\addlinespace[4pt]
\textbf{3. Agency} & \texttt{"dilli police aur sikyoriti phorses ne aam janta par lathi charj kar diya. aansoo gais chhod die. pelet gans ki bauchhar kar di."} \texttt{[03:39--03:47]} vs \texttt{"ladakiyon ko seksual hais karna... kuch logon dvara kiya gaya jo na vardi mein the na nem taig the."} \texttt{[03:57--04:07]} & Places direct active agency (\texttt{ne... kar diya}) on Delhi Police for tear gas and pellet guns, while highlighting how removing uniforms and name-tags deliberately obscures institutional agency (\texttt{kuch logon dvara}) to evade accountability. \\
\addlinespace[4pt]
\textbf{4. Antagonisation} & \texttt{"pavar of pipal ij gretar den da pipal in pavar."} \texttt{[12:39]} \newline  \texttt{"hamara media jisko aapki aavaj uthani thi usi ne aapki aavaj ka daman kiya."} \texttt{[21:56]} & Positions organic, civic-minded youth and ordinary taxpayers (\texttt{pablik}, \texttt{lakhon bachche}) against an arrogant ruling establishment and complicit \texttt{"godi media"} anchors (\texttt{"doosra Goswami"}, \texttt{[22:22]}) who smear students as \texttt{"terarist"} and \texttt{"enti neshanal"}. \\
\addlinespace[4pt]
\textbf{5. Intensification / softening} & \textbf{Intensification:} \texttt{"pelet gans ki bauchhar kar di"} \texttt{[03:39]}, \texttt{"jan sailab aane laga"} [a human flood arrived] \texttt{[03:12]}, \texttt{"enkar ko dalali karne se nahin roka"} [stop anchors from pimping/brokering] \texttt{[23:13]} \newline  \textbf{Softening (protester clashes):} \texttt{"ritalieshan mein un pe bhi hamla huaa tha"} \texttt{[04:45]} & Heightens state repression and media complicity via flood/barrage metaphors and harsh moral vocabulary (\texttt{dalali}), while softening protester violence against 118 injured police personnel by framing it as reactive \texttt{"ritalieshan"}. \\
\bottomrule
\end{longtable}

\subsection*{Short summary (Event 2)}
\begin{itemize}[leftmargin=1.5em, itemsep=3pt, topsep=4pt]
  \item \textbf{What interpretation of the event does this video encourage?} That the NEET/CJP protest was a historic, organic, non-partisan youth movement (\texttt{[00:40]}, \texttt{[15:39]}) triggered by systemic educational rot and paper leaks (\texttt{[15:59]}), which overcame brutal police repression on 20 July and mainstream media defamation to force Union Education Minister Dharmendra Pradhan's resignation (\texttt{[00:08--00:48]}, \texttt{[08:16--08:52]}).
  \item \textbf{Whose perspective is foregrounded, challenged, or missing?} \textit{Foregrounded:} Protesting Gen-Z students, middle/lower-class parents burdened by private school fees (\texttt{[19:09--19:35]}), and ground TV reporters caught between angry crowds and toxic studio anchors (\texttt{[22:07--22:22]}). \textit{Challenged:} Government delay, university gag circulars (\texttt{[15:00--15:20]}), and Republic-style "terrorist/anti-national" smears (\texttt{[05:31--05:56]}, \texttt{[22:22--23:04]}). \textit{Missing:} Direct first-person interviews with NTA/Ministry officials or injured police officers explaining crowd-control failures.
  \item \textbf{What evidence supports your interpretation? What complicates it?} \textit{Supports:} MEA's confirmation of zero foreign funding (\texttt{[00:40--00:48]}), specific GDP education budget trends (\texttt{4.5\%} down to \texttt{2.5\%} at \texttt{[17:52--18:10]}; \texttt{86 lakh} students leaving government schools at \texttt{[19:19]}), and firsthand observation of Jantar Mantar's medical/library/food camps (\texttt{[07:21--08:08]}). \textit{Complicates:} He acknowledges CJP founder Abhijit Dipke's past AAP social-media role (\texttt{[02:19]}), presence of hooting \texttt{"ansoshal eliments"} near CP (\texttt{[07:12--07:21]}), and injuries to 118 police personnel (\texttt{[04:37--04:53]}), showing the event was more complex than pure unprovoked one-way state action.
\end{itemize}
\vspace{4pt}\noindent\rule{\textwidth}{0.4pt}\vspace{4pt}

\section*{Common form: Blank excerpt record}
\textit{Copy or print this page whenever you find a passage worth analysing. Keep the exact wording and enough context to interpret it.}


\subsection*{Excerpt Record 1 (Event 1 --- Ram Mandir Donation Controversy)}
\begin{itemize}[leftmargin=1.5em, itemsep=3pt, topsep=4pt]
  \item \textbf{Outlet:} Sarthak Goswami (\textit{Sunday Show}) | \textbf{Date:} \texttt{2026-07-05}
  \item \textbf{Video title:} \textit{Ram Mandir Scam, Gadkari's Ethanol, Kejriwal's Temple \& More | Sunday Show}
  \item \textbf{Timestamp(s):} \texttt{[00:37--00:56]} | \textbf{Analyst:} Baani Solanki | \textbf{Language:} Hindi / Hinglish (Roman script)
  \item \textbf{Exact transcript excerpt:}
    \begin{quote}\small\ttfamily [00:37] apni report dayar ki jiske baad Trust ke aath logon ko arest kiya gaya. is poore ghotale ka mastarmaind bataya ja raha hai tinnoo e Ram shankar yadav... [00:46] champat raay RSS VHP ka aadmi hai aur tinnoo draivar bali ka bakra. [00:56] tinnoo par aarop hai ki usne apne malik champat raay ki pahunch ka phayda uthakar apne sare phaimili membars ko kaunting tim ka part bana liya.\end{quote}
  \item \textbf{Feature(s):} $\boxtimes$ Naming $\boxtimes$ Presupposition $\boxtimes$ Agency $\square$ Antagonisation $\boxtimes$ Intensification / softening
  \item \textbf{Who is speaking? Is this the outlet's wording or a quotation?} Sarthak Goswami speaking in his own editorial voice, summarizing official SIT claims (\texttt{bataya ja raha hai}) before interjecting his own judgement (\texttt{bali ka bakra}).
  \item \textbf{What does the wording imply, foreground, soften, or leave unstated?} The passive construction \texttt{"mastarmaind bataya ja raha hai"} [is being called the mastermind] distances Sarthak from the official police theory, while the copular assertion \texttt{"champat raay RSS VHP ka aadmi hai aur tinnoo draivar bali ka bakra"} presupposes political protection for Champat Rai and frames the driver as a sacrificial fall guy.
  \item \textbf{What context or evidence supports or complicates this interpretation?} Supported by the stark power asymmetry between a Rs.\,12,000--20,000/month driver and the VHP Vice-President running the Trust (\texttt{[01:14--01:40]}). Complicated by CCTV and asset-recovery reports showing counting staff personally amassed lakhs and crores (\texttt{[01:14--01:31]}), meaning Tinnu was an active participant, not an innocent bystander.
  \item \textbf{How does another outlet describe the same event or actor?} \textbf{Republic TV} (\texttt{0H-bruoEkaY}, \texttt{[03:22--03:42]} \& \texttt{[07:05--07:14]}) describes the exact same 8 arrested men strictly as \texttt{"low-level functionaries in the Ram Mandir trust, some cashiers"} to insulate senior leadership and praise \texttt{"Yogi's zero tolerance"}, whereas \textbf{Aaj Tak} panelist Sangeet Ragi (\texttt{B4L1pz0o4dc}, \texttt{[12:25--12:56]}) defends Champat Rai as an ascetic \texttt{"do dhoti aur do kurte par baithne vala aadmi"}.
  \item \textbf{My interpretation, with reasons:} Sarthak uses colloquial idiom (\texttt{bali ka bakra}) and patron-client naming (\texttt{malik champat raay}) to pierce bureaucratic scapegoating, encouraging viewers to question how a personal driver could pack the counting room with relatives without senior complicity.
\end{itemize}

\subsection*{Excerpt Record 2 (Event 1 --- Ram Mandir Donation Controversy)}
\begin{itemize}[leftmargin=1.5em, itemsep=3pt, topsep=4pt]
  \item \textbf{Outlet:} Sarthak Goswami (\textit{Sunday Show}) | \textbf{Date:} \texttt{2026-07-05}
  \item \textbf{Video title:} \textit{Ram Mandir Scam, Gadkari's Ethanol, Kejriwal's Temple \& More | Sunday Show}
  \item \textbf{Timestamp(s):} \texttt{[02:54--03:28]} | \textbf{Analyst:} Baani Solanki | \textbf{Language:} Hindi / Hinglish (Roman script)
  \item \textbf{Exact transcript excerpt:}
    \begin{quote}\small\ttfamily [02:54] dharma chalane ka bachane ka theka toh tum hi ne diya tha in logon ko. ab jab chori pakdi gaee toh tum hi ko kya slaid [gaslight] kar rahe hain ye? Hindu ka paisa agar kisi Hindu ne le bhi [03:02] liya bhaee tum logon ko kyon mirchi lag rahi hai? kahta hai hinduon ka paisa agar Hindu ne chori kar liya toh galat kya hai? ... [03:12] dekh lo yahi hai inke Hindu rashtra ka trelar... [03:20] main yahi log Hindu khatre mein hai. Hindu khatre mein tumhen ghutti pilae ja rahe the. tum pie ja rahe the.\end{quote}
  \item \textbf{Feature(s):} $\boxtimes$ Naming $\boxtimes$ Presupposition $\boxtimes$ Agency $\boxtimes$ Antagonisation $\boxtimes$ Intensification / softening
  \item \textbf{Who is speaking? Is this the outlet's wording or a quotation?} Sarthak Goswami's commentary embedding a paraphrased quotation (\texttt{kahta hai...}) from pro-Hindutva social media defenders.
  \item \textbf{What does the wording imply, foreground, soften, or leave unstated?} Foregrounds ideological betrayal: \texttt{"dharma... ka theka"} [contract/franchise over religion] implies religion was privatized to political brokers; \texttt{"Hindu rashtra ka trelar"} intensifies a temple theft into a dystopian preview of majoritarian impunity where co-religionist crime is excused.
  \item \textbf{What context or evidence supports or complicates this interpretation?} Supported by real viral rationalizations circulating online (also cited independently by Dhruv Rathee in \texttt{oeH0fOSJuX8} at \texttt{[14:34--14:42]}: \textit{"If a Hindu takes the money of a Hindu, why are you people feeling guilty"}). Complicated by the fact that many Hindu organizations and saints (e.g., Nirmohi Akhada cited on The Wire \texttt{7rUL8fGgRJA} at \texttt{[29:20]}) fiercely condemned the theft rather than excusing it.
  \item \textbf{How does another outlet describe the same event or actor?} \textbf{Republic TV} (\texttt{0H-bruoEkaY}, \texttt{[01:23--01:52]}) inverts the accusation entirely, claiming critics of the scam are \texttt{"people who say let's divide Hindus... people who did not want the Ram Mandir itself"}, whereas \textbf{Dhruv Rathee} (\texttt{oeH0fOSJuX8}, \texttt{[14:24--14:42]}) labels such apologists \texttt{"blind devotees who look at Narendra Modi above Lord Ram"}.
  \item \textbf{My interpretation, with reasons:} By quoting the apologist's own communal logic and extending it to murder (\texttt{"Hindu ka khoon Hindu ne bahaya"}, \texttt{[03:12]}), Sarthak exposes how majoritarian identity politics erodes basic legal and moral accountability.
\end{itemize}

\subsection*{Excerpt Record 3 (Event 1 --- Callback on Temple Appeasement)}
\begin{itemize}[leftmargin=1.5em, itemsep=3pt, topsep=4pt]
  \item \textbf{Outlet:} Sarthak Goswami (\textit{Sunday Show}) | \textbf{Date:} \texttt{2026-07-05}
  \item \textbf{Video title:} \textit{Ram Mandir Scam, Gadkari's Ethanol, Kejriwal's Temple \& More | Sunday Show}
  \item \textbf{Timestamp(s):} \texttt{[27:37--28:13]} | \textbf{Analyst:} Baani Solanki | \textbf{Language:} Hindi / Hinglish (Roman script)
  \item \textbf{Exact transcript excerpt:}
    \begin{quote}\small\ttfamily [27:37] disaid kar lo. dharma ki rajniti rokni hai ya karni hai. kyunki wo aap BJP se betar nahin kar paoge. aap sasti BJP banke rah jaoge... [27:55] BJP ke khilaph jate ho BJP ban jate hain. aise hi wo CJP vala hai. wo abhi risentli karne laga ki ham dilli police ne chhatrapati shivaji maharaj aur bhagat sinh ki kitab kaise phenki maphi mango... [28:04] tum wahan pe ejukeshan ke lie baithe ho rashtravad ke lie. BJP ka gem bajane aate hain, BJP ka hi gem khelne lagte hain.\end{quote}
  \item \textbf{Feature(s):} $\boxtimes$ Naming $\boxtimes$ Presupposition $\square$ Agency $\boxtimes$ Antagonisation $\boxtimes$ Intensification / softening
  \item \textbf{Who is speaking? Is this the outlet's wording or a quotation?} Sarthak Goswami's direct editorial commentary critiquing Arvind Kejriwal/AAP and briefly warning CJP protesters.
  \item \textbf{What does the wording imply, foreground, soften, or leave unstated?} Coining \texttt{"sasti BJP"} [cut-price/cheap BJP] presupposes BJP owns the copyright on religious nationalism (\texttt{BJP ka gem}) and implies that secular/issue-based challengers lose credibility whenever they adopt religious or hyper-nationalist symbolism.
  \item \textbf{What context or evidence supports or complicates this interpretation?} Supported by Punjab AAP's Rs.\,80-crore temple redevelopment and taxpayer-funded pilgrimage announcements (\texttt{[26:14--26:35]}). Complicates any simplistic reading of Sarthak as an uncritical opposition or CJP cheerleader---even 3 weeks before CJP's victory, he critiques CJP for drifting from \texttt{"ejukeshan"} into \texttt{"rashtravad"} symbolism (\texttt{[28:04]}).
  \item \textbf{How does another outlet describe the same event or actor?} \textbf{Republic TV} (\texttt{0H-bruoEkaY}, \texttt{[12:06--12:36]}) also attacks Kejriwal's temple politics as \texttt{"opportunistic"}, though from a pro-BJP angle demanding he \texttt{"atone for one's sins"} before Lord Ram.
  \item \textbf{My interpretation, with reasons:} This passage demonstrates ideological consistency: Sarthak critiques religious instrumentalization across party lines, strengthening his credibility when he defends student protesters in Event 2.
\end{itemize}

\subsection*{Excerpt Record 4 (Event 2 --- NEET Paper Leak Protest)}
\begin{itemize}[leftmargin=1.5em, itemsep=3pt, topsep=4pt]
  \item \textbf{Outlet:} Sarthak Goswami (\textit{Sunday Show}) | \textbf{Date:} \texttt{2026-07-26}
  \item \textbf{Video title:} \textit{Dharmendra Pradhan Resigned, CJP won, Modi on Instagram \& more | Sunday Show}
  \item \textbf{Timestamp(s):} \texttt{[03:02--04:17]} | \textbf{Analyst:} Baani Solanki | \textbf{Language:} Hindi / Hinglish (Roman script)
  \item \textbf{Exact transcript excerpt:}
    \begin{quote}\small\ttfamily [03:02] hangar straik par baithe sonam vangachuk ko dilli police aati hai sivil kloths mein aur unko saphed kapde mein lapet ke embulens mein daal ke wahan se abadakt karke le jati hai. not [03:12] diten abadakt... end den kem da blaik mande 20th julaee 2026... [03:39] dilli police aur sikyoriti phorses ne aam janta par lathi charj kar diya. aansoo gais chhod die. pelet gans ki bauchhar kar di... [03:57] ladakiyon ko seksual hais karna, bachchiyon ko marna yeh sab kuch logon dvara kiya gaya jo na vardi mein the na nem taig the. police ke the toh police aisa kyon karti hai? aur agar police ke nahin the toh wo wahan kya kar rahe the police ke saath?\end{quote}
  \item \textbf{Feature(s):} $\boxtimes$ Naming $\square$ Presupposition $\boxtimes$ Agency $\boxtimes$ Antagonisation $\boxtimes$ Intensification / softening
  \item \textbf{Who is speaking? Is this the outlet's wording or a quotation?} Sarthak Goswami's own narrative reconstruction of the 18 July and 20 July events in Delhi.
  \item \textbf{What does the wording imply, foreground, soften, or leave unstated?} Self-repairing \texttt{"not diten abadakt"} explicitly strips state action of lawful policing legitimacy. The dilemma posed at \texttt{[04:07]} (\texttt{police ke the toh... agar police ke nahin the toh...}) traps the administration in an inescapable accountability bind over unbadged men assaulting \texttt{"ladakiyon"} and \texttt{"bachchiyon"} [young girls].
  \item \textbf{What context or evidence supports or complicates this interpretation?} Corroborated across multiple independent testimonies---Dhruv Rathee (\texttt{l-6Pm-112io}, \texttt{[05:05--05:14]} \& \texttt{[11:37--12:25]}), Mohak Mangal's injured student interviewees (\texttt{v-RikSqb0Vw}, \texttt{[01:38]} \& \texttt{[04:38]}), and Naseeruddin Shah on The Wire (\texttt{pYaTVvHeGoM}, \texttt{[07:34--07:42]}). Complicated by Delhi Police's counter-claim (on Aaj Tak \texttt{JXLIs0WtJAM}, \texttt{[19:43--23:14]}) that officers faced heavy stone-pelting at barricades, which Sarthak himself notes at \texttt{[04:45--04:53]}.
  \item \textbf{How does another outlet describe the same event or actor?} \textbf{Republic TV} (\texttt{HTvb6E\_1Qdc}, \texttt{[01:37--01:47]} \& \texttt{[03:36--04:19]}) completely omits police pellet guns, plainclothes men, and Wangchuk's removal, instead portraying \texttt{"our police, our RAF"} solely as victims attacked by \texttt{"rabid, bloodthirsty mobs"} and \texttt{"goons of Abhijit Dipke"}.
  \item \textbf{My interpretation, with reasons:} The metalinguistic correction (\texttt{"not diten abadakt"}) is textbook critical discourse framing: Sarthak foregrounds how state vocabulary ("detention", "crowd control") sanitizes irregular force.
\end{itemize}

\subsection*{Excerpt Record 5 (Event 2 --- NEET Paper Leak Protest)}
\begin{itemize}[leftmargin=1.5em, itemsep=3pt, topsep=4pt]
  \item \textbf{Outlet:} Sarthak Goswami (\textit{Sunday Show}) | \textbf{Date:} \texttt{2026-07-26}
  \item \textbf{Video title:} \textit{Dharmendra Pradhan Resigned, CJP won, Modi on Instagram \& more | Sunday Show}
  \item \textbf{Timestamp(s):} \texttt{[07:21--08:16]} | \textbf{Analyst:} Baani Solanki | \textbf{Language:} Hindi / Hinglish (Roman script)
  \item \textbf{Exact transcript excerpt:}
    \begin{quote}\small\ttfamily [07:21] par jo apne aap mein protest sait thi. bahut hi achchhi kamyuniti spes ban gaee thi. ek hyooman chen logon ne banaee huee hai taki rasta alag rakha [07:30] ja sake. wahan pe koi log postar bana raha hai, koi khana bant raha hai... paas mein ek medikal kaimp laga huaa hai... [07:44] ek laibreri bani huee hai toh apne log kitaben donet kar rahe hain... [07:58] Zomato Swiggi vale khade hue hain khana le gae. dait ij hau sarvaiv kiya poora moovment. toh ek taim pe jo lidarles moovment tha na uska lidar khud pablik hi ban [08:08] gaee... jab koi nahin tha kisi ke lie toh sab the sabke lie.\end{quote}
  \item \textbf{Feature(s):} $\boxtimes$ Naming $\square$ Presupposition $\boxtimes$ Agency $\square$ Antagonisation $\boxtimes$ Intensification / softening
  \item \textbf{Who is speaking? Is this the outlet's wording or a quotation?} Sarthak Goswami speaking as a firsthand eyewitness who visited Jantar Mantar (\texttt{[08:08]}).
  \item \textbf{What does the wording imply, foreground, soften, or leave unstated?} Renames a contested protest zone a \texttt{"bahut hi achchhi kamyuniti spes"} [a wonderful community space] and distributes constructive civic agency across anonymous citizens (\texttt{koi postar bana raha hai}, \texttt{koi khana bant raha hai}), countering "violent mob" stereotypes.
  \item \textbf{What context or evidence supports or complicates this interpretation?} Supported by ground visuals of medical/food aid and peaceful sit-ins echoed by Mohak Mangal's ground interviewees (\texttt{v-RikSqb0Vw}, \texttt{[05:57]}, \texttt{[10:36--10:47]}). Complicated by Sarthak's own admission seconds earlier (\texttt{[07:12--07:21]}) that adjacent Connaught Place lanes had \texttt{"ansoshal eliments"} arriving in vans and hooting.
  \item \textbf{How does another outlet describe the same event or actor?} \textbf{Republic TV} (\texttt{HTvb6E\_1Qdc}, \texttt{[04:46--05:26]}) characterizes the exact same Jantar Mantar site as a scene of \texttt{"drunken violence"}, \texttt{"sexual perversion"}, and anti-national \texttt{"Nepal model"} anarchy where \texttt{"the students had evaporated"}.
  \item \textbf{My interpretation, with reasons:} By separating the core protest camp from peripheral troublemakers (\texttt{[07:12--07:21]}) and detailing mundane acts of care (libraries, medicine lists, Swiggy orders), Sarthak constructs democratic legitimacy from below.
\end{itemize}

\subsection*{Excerpt Record 6 (Event 2 --- Media Critique \& Anti-National Labelling)}
\begin{itemize}[leftmargin=1.5em, itemsep=3pt, topsep=4pt]
  \item \textbf{Outlet:} Sarthak Goswami (\textit{Sunday Show}) | \textbf{Date:} \texttt{2026-07-26}
  \item \textbf{Video title:} \textit{Dharmendra Pradhan Resigned, CJP won, Modi on Instagram \& more | Sunday Show}
  \item \textbf{Timestamp(s):} \texttt{[22:22--23:04]} | \textbf{Analyst:} Baani Solanki | \textbf{Language:} Hindi / Hinglish (Roman script)
  \item \textbf{Exact transcript excerpt:}
    \begin{quote}\small\ttfamily [22:22] aaj bhi wo doosra Goswami laga hi huaa hai. yeh enti neshanal hain. ye terarist hain... [22:30] da kangrechuleshans CJP nau ophishiyli terarist sar kaming aaut. itne ghatiya, itne daphar, itne dam aargument men strim ho gae na... [22:48] lakhon lakhon bachche sadkon pe dikh rahe hain. sab entineshanal hain. sab pakistan phanded hain. ab stoodiyo mein jane se pahle akhbar padh le. pakistan khud \$10 biliyan dolar ka [22:56] lon mang raha hai amerika se. wo kya phand karega?\end{quote}
  \item \textbf{Feature(s):} $\boxtimes$ Naming $\boxtimes$ Presupposition $\square$ Agency $\boxtimes$ Antagonisation $\boxtimes$ Intensification / softening
  \item \textbf{Who is speaking? Is this the outlet's wording or a quotation?} Sarthak Goswami's wording parodying and directly responding to Republic TV anchor Arnab Goswami (\texttt{"doosra Goswami"}).
  \item \textbf{What does the wording imply, foreground, soften, or leave unstated?} Foregrounds how state-aligned TV channels weaponize national-security labels (\texttt{enti neshanal}, \texttt{terarist}, \texttt{pakistan phanded}) to delegitimize domestic youth grievances, and punctures the conspiracy via geopolitical satire (Pakistan begging for a \$10B bailout).
  \item \textbf{What context or evidence supports or complicates this interpretation?} Directly supported by MEA's official denial of foreign funding (\texttt{[00:40--00:48]}) and by Republic TV's actual broadcast (\texttt{HTvb6E\_1Qdc}, \texttt{[00:23--00:33]} calling Dipke \texttt{"Deep State Dipke"} and \texttt{[04:19--04:46]} contrasting protesters with soldiers fighting terrorists in Anantnag).
  \item \textbf{How does another outlet describe the same event or actor?} \textbf{Aaj Tak} (\texttt{JXLIs0WtJAM}, \texttt{[10:04--11:10]}) avoids calling students terrorists, instead airing Maharashtra CM Devendra Fadnavis's softer delegitimization that outsiders \texttt{"apna-apna ejenda chala rahe hain"} without even knowing what NEET is.
  \item \textbf{My interpretation, with reasons:} This meta-discursive critique directly exemplifies the Language \& Society theme: Sarthak exposes \textit{Naming} ("terrorist" vs \texttt{"lakhon bachche"}) as a deliberate ideological tool of media delegitimization.
\end{itemize}
\vspace{4pt}\noindent\rule{\textwidth}{0.4pt}\vspace{4pt}

\section*{Common form: End-of-event summary}

\subsection*{End-of-Event Summary --- Event 1 (Ram Mandir Donation Controversy)}
\begin{itemize}[leftmargin=1.5em, itemsep=3pt, topsep=4pt]
  \item \textbf{The main way this outlet framed the event was:} Not merely as an Rs.\,80-lakh/crore accounting theft by temple staff, but as a deeper moral and political \texttt{"asli ghotala"} [real scam] resulting from the commercialization of Hindu faith (\texttt{"dharma ka dhandha"}) and citizens outsourcing religious authority to unaccountable political-ideological contractors (\texttt{[00:00--00:56]}, \texttt{[02:12--06:26]}).
  \item \textbf{The strongest linguistic evidence was (quote + timestamp):} \texttt{"champat raay RSS VHP ka aadmi hai aur tinnoo draivar bali ka bakra... Hindu jago Ram ka chanda nahin Ram ka charitra chori ho chuka hai."} (\texttt{[00:46]} \& \texttt{[06:26]}).
  \item \textbf{The shared feature(s) most clearly illustrated were:} \textbf{Naming} (\texttt{bali ka bakra}, \texttt{dharma ka dhandha}, \texttt{charitra chori}), \textbf{Presupposition} (\texttt{tum sarapraizd kyon ho?} at \texttt{[02:12]}), and \textbf{Agency} (holding voters/devotees active agents who handed over the \texttt{"theka"} of religion at \texttt{[02:46--03:38]}).
  \item \textbf{An alternative interpretation of the same evidence is:} That the embezzlement was an operational lapse in cash-counting protocols by colluding counting staff and bank supervisors---not an indictment of all temple devotees or of the Ram Mandir movement itself---and that the UP government's SIT arrest of 8 suspects (\texttt{[00:37]}) represented lawful institutional action rather than a cover-up.
  \item \textbf{What I cannot conclude from this transcript alone:} Whether senior trustees like Champat Rai personally profited from or directed Tinnu Yadav's theft, what the final SIT forensic audit concluded, or the exact rupee total stolen across cash and jewellery.
\end{itemize}

\subsection*{End-of-Event Summary --- Event 2 (Students' Protest Regarding NEET Paper Leak)}
\begin{itemize}[leftmargin=1.5em, itemsep=3pt, topsep=4pt]
  \item \textbf{The main way this outlet framed the event was:} As a victorious, organic, citizen-led democratic awakening (\texttt{"lidar khud pablik hi ban gaee"}, \texttt{[08:08]}) rooted in genuine educational despair (paper leaks, suicides, falling GDP share for education) that defeated police brutality (\texttt{20 July Black Monday}) and \texttt{"godi media"} smear campaigns to compel ministerial accountability (\texttt{[00:08--01:26]}, \texttt{[15:39--18:44]}).
  \item \textbf{The strongest linguistic evidence was (quote + timestamp):} \texttt{"sonam vangachuk ko dilli police aati hai sivil kloths mein... wahan se abadakt karke le jati hai. not diten abadakt... jab koi nahin tha kisi ke lie toh sab the sabke lie."} (\texttt{[03:02--03:12]} \& \texttt{[08:08--08:16]}).
  \item \textbf{The shared feature(s) most clearly illustrated were:} \textbf{Naming} (\texttt{not diten abadakt}, \texttt{blaik mande}, \texttt{kamyuniti spes} vs TV anchors' \texttt{terarist}/\texttt{enti neshanal}), \textbf{Agency} (foregrounding police lathi/pellet-gun force and unbadged plainclothes men at \texttt{[03:39--04:17]}), and \textbf{Antagonisation} (\texttt{pavar of pipal} vs \texttt{pipal in pavar} and \texttt{"dalali"} anchors at \texttt{[12:39]}, \texttt{[23:13]}).
  \item \textbf{An alternative interpretation of the same evidence is:} That once an unauthorized crowd breached Section 163 barricades to march toward Parliament on 20 July and injured 118 police personnel (\texttt{[03:30--03:39]}, \texttt{[04:37--04:53]}), security forces faced an acute law-and-order breakdown, and the government's eventual acceptance of demands reflected negotiated democratic responsiveness rather than mere bruised \texttt{"eego"} (\texttt{[00:40]}).
  \item \textbf{What I cannot conclude from this transcript alone:} The exact identity/command chain of the plainclothes individuals seen without name-tags (\texttt{[04:07]}), how many injuries on 20 July were caused by initial police action versus crowd stone-pelting, or whether post-resignation structural NTA reforms actually prevented future leaks.
\end{itemize}
\textit{Research reminder: interpret wording in context, distinguish allegations from established facts, and support claims with exact excerpts. Mixed or contradictory evidence is worth recording.}


\end{document}
